\documentclass[12pt]{letter}

% \usepackage[britdate]{canterbury-letter}
\usepackage[britdate]{LiU-letter}
\usepackage{times}
\usepackage{letterbib}
\usepackage{geometry}
\usepackage[round]{natbib}
\usepackage{graphicx}
\geometry{a4paper}
\usepackage[T1]{fontenc}
\usepackage[utf8]{inputenc}
\usepackage{authblk}
\usepackage[running]{lineno}
\usepackage{amsmath,amsfonts,amssymb}
% \usepackage[margin=10pt,font=small,labelfont=bf]{caption}

%\usepackage{natbib}
% \bibpunct[; ]{(}{)}{;}{a}{,}{;}

\newenvironment{refquote}{\bigskip \begin{it}}{\end{it}\smallskip}

\newenvironment{figure}{}

\position{Associate professor}
\department{Department of Physics, Chemistry, and Biology (IFM)}
\location{581 83 Link\"{o}ping, Sweden}
\cityzip{581 83 Link\"{o}ping, Sweden}
\telephone{}
\fax{}
\email{anna.eklof@liu.se}
\url{https://liu.se/en/employee/annek49}
\name{ Anna C Ekl\"{o}f}


\newcommand{\mytitle}{\emph{Conditions for successful range shifts under climate change -- the role of species dispersal and landscape configuration }}
\newcommand{\myjournal}{\emph{Diversity and Distributions}}

\begin{document}

\begin{letter}{\bf Distinguished  Professor Janet Franklin\\
               Editor-in-Chief, Diversity And Distributions\\
               Department of Botany and Plant Sciences\\
               University of California - Riverside\\
               900 University Ave\\
               Riverside CA 92521, USA
                               }


\opening{Dear Prof. Franklin:}
My co-authors and I are pleased to submit our manuscript \mytitle for publication in \myjournal. 
    
Ongoing climate change force many species around the globe to modify their current geographical distributions. Several factors are of importance for if the species will be able to successfully move to and colonize new areas, i.e. successfully shift their range. Predicting biodiversity scenarios under climate change has so far largely relied on species distribution models, although the accuracy of these methods have recently been highly debated as they do not account for, for example, species dispersal abilities and population dynamics. 

Here we develop a novel theoretical framework for modeling the range shift of a single species where we specifically model species population dynamics and dispersal ability. We have chosen to focus on three factors recognized for their importance for species range shifts; {\em i)} dispersal distance for the species, {\em ii)}  the configuration for suitable habitat in the regional landscape, and {\em iii)} speed of climate change, and analyze the effect of these factors in combined. To justify the ecological relevance of our model, i.e., model parameter values, we use empirical data on species dispersal distances and current and future landscape configuration for a set of European mammals. 

We show that all factors are of importance for the outcomes of species range shifts, but in particular the specific configuration of suitable habitat in the landscape have far reaching consequences. In landscapes with a homogeneous distribution of habitats, all but the least 
dispersive species will be able to range shift. In landscapes were the suitable habitats are distributed more heterogeneously species’ ability to range shift is far more variable and heavily dependent on the actual habitat configuration. Our results highlight the importance of taking the interplay between species dispersal, speed of climate change and landscape configuration into account when predicting future species distributions in the face of climate change. The results also have implications for conservations planning such as the highly actual organization of green infrastructure elements.

    Because of this, we hope that you will agree that our manuscript merits 
    publication in \myjournal.  We look forward to your response.


\closing{Best regards,}


\end{letter}



\end{document}
